\documentclass[a4paper, 10pt]{article}

\usepackage{amsmath}
\usepackage{amssymb}
\usepackage{booktabs}
\usepackage{geometry}
\usepackage{kotex}
\usepackage{float}

\begin{document}

    \section{$Y$ per capita}
        이전 미팅자료의 section 3과 동일하게 세 가지 접근성 변수를 사용하되,
        종속변수를 인구 만 명당 outcome으로 바꾸어 추정하였다.
        병원 분석에서 $hospital\_no$, $net\_entry$, $exit\_no$,
        $entry\_no$는 인구 만 명당 값으로 변환하였다. $exit\_rate$와
        $entry\_rate$는 이미 비율 변수이므로 원래 단위를 유지하였다. 의사 분석에서는
        각 doctor level별 의사 수를 인구 만 명당 의사 수로 변환하였다.

        모든 회귀식은 서울, 경기, 인천, 제주를 제외한 2007--2025년 시군구-연도 panel을
        사용한다. 이전 분석과 달리 종속변수가 population으로 normalize되어 있으므로
        $\log(Population_{it})$ control은 포함하지 않는다. 모든 specification에는
        region fixed effects와 year fixed effects를 포함하고, standard error는 region
        level에서 cluster하였다.

        \subsection{Travel time}
            첫 번째 분석은 terminal까지의 절대적 접근시간과 인구 만 명당 의료공급 사이의
            관계를 추정한다.
            \begin{align*}
                Y^{pc}_{it}
                =
                \alpha_i + \tau_t
                + \beta TravelTime_{it}
                + \epsilon_{it}
            \end{align*}

            \begin{table}[H]
                \centering
                \caption{Per-Capita Outcome, Travel Time: 의원}
                \begingroup
\centering
\begin{tabular}{lcccccc}
   \toprule
                          & hospital\_no  & net\_entry  & exit\_no  & exit\_rate  & entry\_no  & entry\_rate \\    
                          & (1)           & (2)         & (3)       & (4)         & (5)        & (6)\\  
   \midrule 
   $TravelTime_{it}$      & 0.1456$^{*}$  & 0.0100      & 0.0026    & -0.0011     & 0.0112     & 0.0004\\   
                          & (0.0857)      & (0.0143)    & (0.0138)  & (0.0028)    & (0.0177)   & (0.0039)\\   
   Log population control & No            & No          & No        & No          & No         & No\\  
   Region FE              & Yes           & Yes         & Yes       & Yes         & Yes        & Yes\\  
   Year FE                & Yes           & Yes         & Yes       & Yes         & Yes        & Yes\\  
   Regions                & 159           & 159         & 159       & 159         & 159        & 159\\  
    \\
   Observations           & 3,021         & 2,862       & 3,021     & 3,021       & 2,862      & 2,862\\  
   Adjusted R$^2$         & 0.96404       & 0.11614     & 0.34126   & 0.14672     & 0.26067    & 0.06264\\  
   \bottomrule
\end{tabular}
\par\endgroup

            \end{table}

            \begin{table}[H]
                \centering
                \caption{Per-Capita Outcome, Travel Time: 병원}
                \begingroup
\centering
\begin{tabular}{lcccccc}
   \toprule
                          & hospital\_no  & net\_entry  & exit\_no       & exit\_rate  & entry\_no     & entry\_rate \\    
                          & (1)           & (2)         & (3)            & (4)         & (5)           & (6)\\  
   \midrule 
   $TravelTime_{it}$      & -0.0236       & -0.0017     & -0.0096$^{**}$ & -0.0188     & -0.0114$^{*}$ & 0.0080\\   
                          & (0.0212)      & (0.0055)    & (0.0045)       & (0.0155)    & (0.0059)      & (0.0232)\\   
   Log population control & No            & No          & No             & No          & No            & No\\  
   Region FE              & Yes           & Yes         & Yes            & Yes         & Yes           & Yes\\  
   Year FE                & Yes           & Yes         & Yes            & Yes         & Yes           & Yes\\  
   Regions                & 151           & 151         & 151            & 151         & 151           & 151\\  
    \\
   Observations           & 2,869         & 2,718       & 2,869          & 2,730       & 2,718         & 2,587\\  
   Adjusted R$^2$         & 0.72338       & 0.00785     & 0.08963        & 0.04795     & 0.05788       & 0.05994\\  
   \bottomrule
\end{tabular}
\par\endgroup

            \end{table}

            \begin{table}[H]
                \centering
                \caption{Per-Capita Outcome, Travel Time: 종합병원}
                \input{tables/pc_travel_time_hospital_general.tex}
            \end{table}

            \begin{table}[H]
                \centering
                \caption{Per-Capita Outcome, Travel Time: 상급종합병원}
                \begingroup
\centering
\begin{tabular}{lcccccc}
   \toprule
                          & hospital\_no  & net\_entry  & exit\_no                & exit\_rate  & entry\_no  & entry\_rate \\    
                          & (1)           & (2)         & (3)                     & (4)         & (5)        & (6)\\  
   \midrule 
   $TravelTime_{it}$      & -0.0303$^{*}$ & 0.0046      & $-9.13\times 10^{-5}$   & -0.0074     & 0.0045     & 0.1821\\   
                          & (0.0164)      & (0.0052)    & ($9.58\times 10^{-5}$)  & (0.0083)    & (0.0052)   & (0.1390)\\   
   Log population control & No            & No          & No                      & No          & No         & No\\  
   Region FE              & Yes           & Yes         & Yes                     & Yes         & Yes        & Yes\\  
   Year FE                & Yes           & Yes         & Yes                     & Yes         & Yes        & Yes\\  
   Regions                & 20            & 20          & 20                      & 20          & 20         & 20\\  
    \\
   Observations           & 380           & 360         & 380                     & 279         & 360        & 266\\  
   Adjusted R$^2$         & 0.89783       & -0.03386    & -0.00118                & -0.01386    & -0.03439   & 0.06863\\  
   \bottomrule
\end{tabular}
\par\endgroup

            \end{table}

            \begin{table}[H]
                \centering
                \caption{Per-Capita Outcome, Travel Time: 의사}
                \begingroup
\centering
\begin{tabular}{lcccc}
   \toprule
                          & 인턴     & 레지던트 & 전문의   & 일반의 \\   
                          & (1)      & (2)      & (3)      & (4)\\  
   \midrule 
   $TravelTime_{it}$      & -0.0054  & -0.1317  & -0.1594  & -0.0346\\   
                          & (0.0589) & (0.1677) & (0.4706) & (0.0654)\\   
   Log population control & No       & No       & No       & No\\  
   Region FE              & Yes      & Yes      & Yes      & Yes\\  
   Year FE                & Yes      & Yes      & Yes      & Yes\\  
   Regions                & 160      & 160      & 160      & 160\\  
    \\
   Observations           & 3,040    & 3,040    & 3,040    & 3,040\\  
   Adjusted R$^2$         & 0.82300  & 0.90115  & 0.94894  & 0.78557\\  
   \bottomrule
\end{tabular}
\par\endgroup

            \end{table}

        \subsection{Travel-time saving}
            두 번째 분석은 2007년 대비 terminal 접근시간 단축폭과 인구 만 명당 의료공급
            사이의 관계를 추정한다.
            \begin{align*}
                Y^{pc}_{it}
                =
                \alpha_i + \tau_t
                + \beta Saving_{it}
                + \epsilon_{it}
            \end{align*}

            \begin{table}[H]
                \centering
                \caption{Per-Capita Outcome, Travel-Time Saving: 의원}
                \input{tables/pc_saving_hospital_clinic.tex}
            \end{table}

            \begin{table}[H]
                \centering
                \caption{Per-Capita Outcome, Travel-Time Saving: 병원}
                \input{tables/pc_saving_hospital_secondary.tex}
            \end{table}

            \begin{table}[H]
                \centering
                \caption{Per-Capita Outcome, Travel-Time Saving: 종합병원}
                \begingroup
\centering
\begin{tabular}{lcccccc}
   \toprule
                          & hospital\_no  & net\_entry  & exit\_no      & exit\_rate  & entry\_no  & entry\_rate \\    
                          & (1)           & (2)         & (3)           & (4)         & (5)        & (6)\\  
   \midrule 
   $Saving_{it}$          & 0.0028        & -0.0005     & -0.0023$^{*}$ & -0.0172     & -0.0027    & -0.0187\\   
                          & (0.0061)      & (0.0016)    & (0.0014)      & (0.0108)    & (0.0022)   & (0.0205)\\   
   Log population control & No            & No          & No            & No          & No         & No\\  
   Region FE              & Yes           & Yes         & Yes           & Yes         & Yes        & Yes\\  
   Year FE                & Yes           & Yes         & Yes           & Yes         & Yes        & Yes\\  
   Regions                & 99            & 99          & 99            & 99          & 99         & 99\\  
    \\
   Observations           & 1,881         & 1,782       & 1,881         & 1,678       & 1,782      & 1,594\\  
   Adjusted R$^2$         & 0.80216       & -0.00519    & 0.02492       & 0.05975     & 0.00518    & 0.05291\\  
   \bottomrule
\end{tabular}
\par\endgroup

            \end{table}

            \begin{table}[H]
                \centering
                \caption{Per-Capita Outcome, Travel-Time Saving: 상급종합병원}
                \input{tables/pc_saving_hospital_tertiary.tex}
            \end{table}

            \begin{table}[H]
                \centering
                \caption{Per-Capita Outcome, Travel-Time Saving: 의사}
                \input{tables/pc_saving_doctor.tex}
            \end{table}

        \subsection{KTX within two hours}
            세 번째 분석은 KTX를 이용한 terminal 접근시간이 2시간 이하이고 자동차보다
            빠른 경우를 discrete treatment로 정의하여 추정한다.
            \begin{align*}
                Y^{pc}_{it}
                =
                \alpha_i + \tau_t
                + \delta KTXWithin2h_{it}
                + \epsilon_{it}
            \end{align*}

            \begin{table}[H]
                \centering
                \caption{Per-Capita Outcome, KTX Within Two Hours: 의원}
                \begingroup
\centering
\begin{tabular}{lcccccc}
   \toprule
                          & hospital\_no  & net\_entry  & exit\_no  & exit\_rate  & entry\_no  & entry\_rate \\    
                          & (1)           & (2)         & (3)       & (4)         & (5)        & (6)\\  
   \midrule 
   $KTXWithin2h_{it}$     & -0.1310       & 0.0051      & 0.0129    & 0.0022      & 0.0052     & -0.0015\\   
                          & (0.1034)      & (0.0223)    & (0.0307)  & (0.0049)    & (0.0451)   & (0.0086)\\   
   Log population control & No            & No          & No        & No          & No         & No\\  
   Region FE              & Yes           & Yes         & Yes       & Yes         & Yes        & Yes\\  
   Year FE                & Yes           & Yes         & Yes       & Yes         & Yes        & Yes\\  
   Regions                & 159           & 159         & 159       & 159         & 159        & 159\\  
    \\
   Observations           & 3,021         & 2,862       & 3,021     & 3,021       & 2,862      & 2,862\\  
   Adjusted R$^2$         & 0.96391       & 0.11605     & 0.34129   & 0.14671     & 0.26061    & 0.06265\\  
   \bottomrule
\end{tabular}
\par\endgroup

            \end{table}

            \begin{table}[H]
                \centering
                \caption{Per-Capita Outcome, KTX Within Two Hours: 병원}
                \begingroup
\centering
\begin{tabular}{lcccccc}
   \toprule
                          & hospital\_no  & net\_entry  & exit\_no  & exit\_rate  & entry\_no  & entry\_rate \\    
                          & (1)           & (2)         & (3)       & (4)         & (5)        & (6)\\  
   \midrule 
   $KTXWithin2h_{it}$     & 0.0803$^{**}$ & 0.0020      & 0.0022    & 0.0152      & 0.0021     & -0.0102\\   
                          & (0.0312)      & (0.0100)    & (0.0077)  & (0.0235)    & (0.0071)   & (0.0535)\\   
   Log population control & No            & No          & No        & No          & No         & No\\  
   Region FE              & Yes           & Yes         & Yes       & Yes         & Yes        & Yes\\  
   Year FE                & Yes           & Yes         & Yes       & Yes         & Yes        & Yes\\  
   Regions                & 151           & 151         & 151       & 151         & 151        & 151\\  
    \\
   Observations           & 2,869         & 2,718       & 2,869     & 2,730       & 2,718      & 2,587\\  
   Adjusted R$^2$         & 0.72413       & 0.00783     & 0.08872   & 0.04765     & 0.05717    & 0.05991\\  
   \bottomrule
\end{tabular}
\par\endgroup

            \end{table}

            \begin{table}[H]
                \centering
                \caption{Per-Capita Outcome, KTX Within Two Hours: 종합병원}
                \begingroup
\centering
\begin{tabular}{lcccccc}
   \toprule
                          & hospital\_no  & net\_entry  & exit\_no  & exit\_rate  & entry\_no  & entry\_rate \\    
                          & (1)           & (2)         & (3)       & (4)         & (5)        & (6)\\  
   \midrule 
   $KTXWithin2h_{it}$     & 0.0092        & -0.0006     & -0.0058   & -0.0367     & -0.0068    & -0.0299\\   
                          & (0.0081)      & (0.0009)    & (0.0046)  & (0.0273)    & (0.0049)   & (0.0305)\\   
   Log population control & No            & No          & No        & No          & No         & No\\  
   Region FE              & Yes           & Yes         & Yes       & Yes         & Yes        & Yes\\  
   Year FE                & Yes           & Yes         & Yes       & Yes         & Yes        & Yes\\  
   Regions                & 99            & 99          & 99        & 99          & 99         & 99\\  
    \\
   Observations           & 1,881         & 1,782       & 1,881     & 1,678       & 1,782      & 1,594\\  
   Adjusted R$^2$         & 0.80219       & -0.00523    & 0.02467   & 0.05929     & 0.00509    & 0.05253\\  
   \bottomrule
\end{tabular}
\par\endgroup

            \end{table}

            \begin{table}[H]
                \centering
                \caption{Per-Capita Outcome, KTX Within Two Hours: 상급종합병원}
                \begingroup
\centering
\begin{tabular}{lcccccc}
   \toprule
                          & hospital\_no  & net\_entry  & exit\_no                 & exit\_rate              & entry\_no      & entry\_rate \\    
                          & (1)           & (2)         & (3)                      & (4)                     & (5)            & (6)\\  
   \midrule 
   $KTXWithin2h_{it}$     & 0.0107        & 0.0011      & 0.0031$^{***}$           & 0.1667$^{***}$          & 0.0042$^{***}$ & 0.2450$^{***}$\\   
                          & (0.0085)      & (0.0008)    & ($7.02\times 10^{-19}$)  & ($3.44\times 10^{-9}$)  & (0.0008)       & (0.0232)\\   
   Log population control & No            & No          & No                       & No                      & No             & No\\  
   Region FE              & Yes           & Yes         & Yes                      & Yes                     & Yes            & Yes\\  
   Year FE                & Yes           & Yes         & Yes                      & Yes                     & Yes            & Yes\\  
   Regions                & 20            & 20          & 20                       & 20                      & 20             & 20\\  
    \\
   Observations           & 380           & 360         & 380                      & 279                     & 360            & 266\\  
   Adjusted R$^2$         & 0.89080       & -0.03806    & 0.29423                  & 0.28462                 & -0.03779       & 0.06861\\  
   \bottomrule
\end{tabular}
\par\endgroup

            \end{table}

            \begin{table}[H]
                \centering
                \caption{Per-Capita Outcome, KTX Within Two Hours: 의사}
                \begingroup
\centering
\begin{tabular}{lcccc}
   \toprule
                          & 인턴     & 레지던트 & 전문의   & 일반의 \\   
                          & (1)      & (2)      & (3)      & (4)\\  
   \midrule 
   $KTXWithin2h_{it}$     & 0.0289   & 0.0838   & 0.2743   & 0.0860\\   
                          & (0.0283) & (0.1471) & (0.9177) & (0.1727)\\   
   Log population control & No       & No       & No       & No\\  
   Region FE              & Yes      & Yes      & Yes      & Yes\\  
   Year FE                & Yes      & Yes      & Yes      & Yes\\  
   Regions                & 160      & 160      & 160      & 160\\  
    \\
   Observations           & 3,040    & 3,040    & 3,040    & 3,040\\  
   Adjusted R$^2$         & 0.82300  & 0.90112  & 0.94893  & 0.78558\\  
   \bottomrule
\end{tabular}
\par\endgroup

            \end{table}

    \section{Difference From 2007}
        이 섹션에서는 2007년 대비 변화분을 종속변수로 사용한다. 분석 대상 outcome은
        stock 변수로 제한하였다. 병원 분석에서는 의원, 병원, 종합병원, 상급종합병원의
        병원 수 변화를 사용하고, 의사 분석에서는 인턴, 레지던트, 전문의, 일반의 수의
        변화를 사용한다. 독립변수는 2007년 대비 terminal 접근시간 단축폭인
        $Saving_{it}$만 사용한다.

        population control 역시 2007년 대비 변화분으로 넣는다. 즉, 회귀식에는
        $\log(Population_{it}) - \log(Population_{i,2007})$가 포함된다. 모든 회귀식은
        region fixed effects와 year fixed effects를 포함하고, standard error는 region
        level에서 cluster하였다.

        병원 stock 변수에 대한 추정식은 다음과 같다.
        \begin{align*}
            HospitalNo^m_{it} - HospitalNo^m_{i,2007}
            =
            \alpha_i + \tau_t
            + \beta Saving_{it}
            + \gamma \left[
                \log(Population_{it}) - \log(Population_{i,2007})
            \right]
            + \epsilon_{it}
        \end{align*}
        여기서 $m$은 의원, 병원, 종합병원, 상급종합병원을 의미한다.

        \begin{table}[H]
            \centering
            \caption{Difference From 2007, Travel-Time Saving: 병원 수}
            \begingroup
\centering
\begin{tabular}{lcccc}
   \toprule
                                     & 의원    & 병원     & 종합병원 & 상급종합병원 \\   
                                     & (1)     & (2)      & (3)      & (4)\\  
   \midrule 
   $Saving_{it}$                     & -1.651  & 0.3255   & 0.0374   & 0.3390$^{**}$\\   
                                     & (2.143) & (0.3448) & (0.0828) & (0.1403)\\   
   Log population difference control & Yes     & Yes      & Yes      & Yes\\  
   Region FE                         & Yes     & Yes      & Yes      & Yes\\  
   Year FE                           & Yes     & Yes      & Yes      & Yes\\  
   Regions                           & 159     & 151      & 99       & 20\\  
    \\
   Observations                      & 3,021   & 2,869    & 1,881    & 380\\  
   Adjusted R$^2$                    & 0.74034 & 0.75648  & 0.68372  & 0.49114\\  
   \bottomrule
\end{tabular}
\par\endgroup

        \end{table}

        의사 stock 변수에 대한 추정식은 다음과 같다.
        \begin{align*}
            DoctorNo^d_{it} - DoctorNo^d_{i,2007}
            =
            \alpha_i + \tau_t
            + \beta Saving_{it}
            + \gamma \left[
                \log(Population_{it}) - \log(Population_{i,2007})
            \right]
            + \epsilon_{it}
        \end{align*}
        여기서 $d$는 인턴, 레지던트, 전문의, 일반의를 의미한다.

        \begin{table}[H]
            \centering
            \caption{Difference From 2007, Travel-Time Saving: 의사 수}
            \begingroup
\centering
\begin{tabular}{lcccc}
   \toprule
                                     & 인턴     & 레지던트 & 전문의  & 일반의 \\   
                                     & (1)      & (2)      & (3)     & (4)\\  
   \midrule 
   $Saving_{it}$                     & 0.0166   & 1.137    & 3.507   & 0.0749\\   
                                     & (0.5804) & (2.208)  & (10.68) & (0.8055)\\   
   Log population difference control & Yes      & Yes      & Yes     & Yes\\  
   Region FE                         & Yes      & Yes      & Yes     & Yes\\  
   Year FE                           & Yes      & Yes      & Yes     & Yes\\  
   Regions                           & 160      & 160      & 160     & 160\\  
    \\
   Observations                      & 3,040    & 3,040    & 3,040   & 3,040\\  
   Adjusted R$^2$                    & 0.82076  & 0.54167  & 0.77020 & 0.50004\\  
   \bottomrule
\end{tabular}
\par\endgroup

        \end{table}

    \section{Population as Outcome}
        이 섹션에서는 의료공급 변수가 아니라 지역 인구 자체를 종속변수로 사용한다.
        구체적으로 $\log(Population_{it})$을 outcome으로 두고, 앞선 분석에서 사용한
        세 가지 접근성 변수인 $TravelTime_{it}$, $Saving_{it}$,
        $KTXWithin2h_{it}$를 각각 독립변수로 넣어 추정하였다. 이 분석에서는
        $\log(Population_{it})$이 종속변수이므로 population control은 포함하지 않는다.

        추정식은 다음과 같다.
        \begin{align*}
            \log(Population_{it})
            =
            \alpha_i + \tau_t
            + \beta Access_{it}
            + \epsilon_{it}
        \end{align*}
        여기서 $Access_{it}$는 $TravelTime_{it}$, $Saving_{it}$,
        $KTXWithin2h_{it}$ 중 하나를 의미한다. 모든 회귀식은 서울, 경기, 인천, 제주를
        제외한 2007--2025년 시군구-연도 panel을 사용하고, region fixed effects와
        year fixed effects를 포함하며, standard error는 region level에서 cluster하였다.

        \begin{table}[H]
            \centering
            \caption{Population as Outcome}
            \begingroup
\centering
\begin{tabular}{lccc}
   \toprule
                      & Travel Time & Saving   & KTX Within 2h \\   
                      & (1)         & (2)      & (3)\\  
   \midrule 
   $TravelTime_{it}$  & 0.0107      &          &   \\   
                      & (0.0179)    &          &   \\   
   $Saving_{it}$      &             & -0.0107  &   \\   
                      &             & (0.0179) &   \\   
   $KTXWithin2h_{it}$ &             &          & -0.0824$^{***}$\\   
                      &             &          & (0.0280)\\   
   Region FE          & Yes         & Yes      & Yes\\  
   Year FE            & Yes         & Yes      & Yes\\  
   Regions            & 160         & 160      & 160\\  
    \\
   Observations       & 3,040       & 3,040    & 3,040\\  
   Adjusted R$^2$     & 0.99380     & 0.99380  & 0.99387\\  
   \bottomrule
\end{tabular}
\par\endgroup

        \end{table}

        추가로 $KTXWithin2h_{it}$에서 관측된 계수가 특정 threshold 선택에 민감한지 확인하기
        위해 threshold를 1시간, 1.5시간, 2시간, 2.5시간, 3시간, 3.5시간, 4시간으로
        바꾸어 동일한 회귀식을 추정하였다. 각 threshold 변수는 KTX를 이용한 terminal
        접근시간이 해당 threshold 이하이고, 동시에 자동차보다 KTX가 빠른 경우 1을 갖는다.
        1시간 threshold는 region fixed effects와 year fixed effects에 의해 완전히 설명되어
        별도의 계수를 추정할 수 없었다.

        \begin{table}[H]
            \centering
            \caption{Population as Outcome, KTX Within Threshold Robustness}
            \begingroup
\centering
\begin{tabular}{lcccccc}
   \toprule
                        & 1.5h            & 2h              & 2.5h          & 3h       & 3.5h     & 4h \\   
                        & (1)             & (2)             & (3)           & (4)      & (5)      & (6)\\  
   \midrule 
   $KTXWithin1.5h_{it}$ & -0.1263$^{***}$ &                 &               &          &          &   \\   
                        & (0.0249)        &                 &               &          &          &   \\   
   $KTXWithin2h_{it}$   &                 & -0.0824$^{***}$ &               &          &          &   \\   
                        &                 & (0.0280)        &               &          &          &   \\   
   $KTXWithin2.5h_{it}$ &                 &                 & -0.0317$^{*}$ &          &          &   \\   
                        &                 &                 & (0.0169)      &          &          &   \\   
   $KTXWithin3h_{it}$   &                 &                 &               & -0.0133  &          &   \\   
                        &                 &                 &               & (0.0169) &          &   \\   
   $KTXWithin3.5h_{it}$ &                 &                 &               &          & -0.0071  &   \\   
                        &                 &                 &               &          & (0.0142) &   \\   
   $KTXWithin4h_{it}$   &                 &                 &               &          &          & -0.0014\\   
                        &                 &                 &               &          &          & (0.0180)\\   
   Region FE            & Yes             & Yes             & Yes           & Yes      & Yes      & Yes\\  
   Year FE              & Yes             & Yes             & Yes           & Yes      & Yes      & Yes\\  
   Regions              & 160             & 160             & 160           & 160      & 160      & 160\\  
    \\
   Observations         & 3,040           & 3,040           & 3,040         & 3,040    & 3,040    & 3,040\\  
   Adjusted R$^2$       & 0.99388         & 0.99387         & 0.99383       & 0.99381  & 0.99380  & 0.99380\\  
   \bottomrule
\end{tabular}
\par\endgroup

        \end{table}

    \appendix

    \section{Hospital Regressions Using All Regions}
        이 appendix에서는 병원 회귀식에 대해, 해당 병원 type의 병원 수가 모든 연도에서 0인
        지역도 제외하지 않고 서울, 경기, 인천, 제주를 제외한 160개 지역 전체를 사용하여
        동일한 분석을 반복한다. 의사 회귀식은 본문 분석에서도 이미 160개 지역 전체를 사용하고
        있으므로 여기서는 병원 회귀식만 보고한다. 단, $EntryRate_{it}$와 $ExitRate_{it}$는
        원자료에서 rate가 정의된 관측치만 사용되므로 표의 region 수가 160보다 작을 수 있다.

        \subsection{Original Section 3 Specifications}
            이 subsection은 이전 미팅자료 section 3의 세 가지 접근성 변수
            $TravelTime_{it}$, $Saving_{it}$, $KTXWithin2h_{it}$를 독립변수로 사용한
            병원 회귀식이다. 종속변수는 병원 수, 순진입, 폐업 수, 폐업률, 진입 수, 진입률이며,
            병원 수는 $\log(HospitalNo_{it}+1)$로 정의하였다. 모든 회귀식은
            $\log(Population_{it})$ control, region fixed effects, year fixed effects를
            포함하고 standard error는 region level에서 cluster하였다.

            \begin{align*}
                Y^m_{it}
                =
                \alpha_i + \tau_t
                + \beta Access_{it}
                + \gamma \log(Population_{it})
                + \epsilon_{it}
            \end{align*}
            여기서 $m$은 의원, 병원, 종합병원, 상급종합병원을 의미하고,
            $Access_{it}$는 $TravelTime_{it}$, $Saving_{it}$, $KTXWithin2h_{it}$ 중 하나이다.

            \subsubsection{Travel Time}
                \begin{table}[H]
                    \centering
                    \caption{All-Region Hospital Regression, Travel Time: 의원}
                    \input{tables/appendix_all_regions_travel_time_hospital_clinic.tex}
                \end{table}

                \begin{table}[H]
                    \centering
                    \caption{All-Region Hospital Regression, Travel Time: 병원}
                    \input{tables/appendix_all_regions_travel_time_hospital_secondary.tex}
                \end{table}

                \begin{table}[H]
                    \centering
                    \caption{All-Region Hospital Regression, Travel Time: 종합병원}
                    \begingroup
\centering
\begin{tabular}{lcccccc}
   \toprule
                            & hospital\_no  & net\_entry  & exit\_no  & exit\_rate  & entry\_no  & entry\_rate \\    
                            & (1)           & (2)         & (3)       & (4)         & (5)        & (6)\\  
   \midrule 
   $TravelTime_{it}$        & -0.0076       & 0.0163      & 0.0250    & 0.0172      & 0.0318     & 0.0192\\   
                            & (0.0285)      & (0.0200)    & (0.0171)  & (0.0108)    & (0.0238)   & (0.0205)\\   
   $\log(Population_{it})$  & 0.3393$^{**}$ & -0.0180     & -0.0027   & 0.0126      & 0.0078     & -0.1442\\   
                            & (0.1451)      & (0.0680)    & (0.0302)  & (0.0445)    & (0.0874)   & (0.1104)\\   
   Log population control   & Yes           & Yes         & Yes       & Yes         & Yes        & Yes\\  
   Region FE                & Yes           & Yes         & Yes       & Yes         & Yes        & Yes\\  
   Year FE                  & Yes           & Yes         & Yes       & Yes         & Yes        & Yes\\  
   Regions                  & 160           & 160         & 160       & 99          & 160        & 99\\  
    \\
   Observations             & 3,040         & 2,880       & 3,040     & 1,678       & 2,880      & 1,594\\  
   Adjusted R$^2$           & 0.92593       & -0.00725    & 0.10498   & 0.05918     & 0.06472    & 0.05436\\  
   \bottomrule
\end{tabular}
\par\endgroup

                \end{table}

                \begin{table}[H]
                    \centering
                    \caption{All-Region Hospital Regression, Travel Time: 상급종합병원}
                    \input{tables/appendix_all_regions_travel_time_hospital_tertiary.tex}
                \end{table}

            \subsubsection{Travel-Time Saving}
                \begin{table}[H]
                    \centering
                    \caption{All-Region Hospital Regression, Travel-Time Saving: 의원}
                    \begingroup
\centering
\begin{tabular}{lcccccc}
   \toprule
                            & hospital\_no   & net\_entry    & exit\_no  & exit\_rate  & entry\_no     & entry\_rate \\    
                            & (1)            & (2)           & (3)       & (4)         & (5)           & (6)\\  
   \midrule 
   $Saving_{it}$            & -0.0146        & -0.3878$^{*}$ & -0.1221   & 0.0013      & -0.3893       & -0.0002\\   
                            & (0.0157)       & (0.2135)      & (0.2222)  & (0.0028)    & (0.2982)      & (0.0040)\\   
   $\log(Population_{it})$  & 0.9241$^{***}$ & 1.113         & 0.6541    & 0.0130      & 2.940$^{***}$ & 0.0215$^{*}$\\   
                            & (0.1544)       & (0.8838)      & (1.198)   & (0.0081)    & (0.7425)      & (0.0129)\\   
   Log population control   & Yes            & Yes           & Yes       & Yes         & Yes           & Yes\\  
   Region FE                & Yes            & Yes           & Yes       & Yes         & Yes           & Yes\\  
   Year FE                  & Yes            & Yes           & Yes       & Yes         & Yes           & Yes\\  
   Regions                  & 160            & 160           & 160       & 159         & 160           & 159\\  
    \\
   Observations             & 3,040          & 2,880         & 3,040     & 3,021       & 2,880         & 2,862\\  
   Adjusted R$^2$           & 0.99452        & 0.33450       & 0.76235   & 0.14685     & 0.71013       & 0.06275\\  
   \bottomrule
\end{tabular}
\par\endgroup

                \end{table}

                \begin{table}[H]
                    \centering
                    \caption{All-Region Hospital Regression, Travel-Time Saving: 병원}
                    \begingroup
\centering
\begin{tabular}{lcccccc}
   \toprule
                            & hospital\_no  & net\_entry  & exit\_no     & exit\_rate  & entry\_no  & entry\_rate \\    
                            & (1)           & (2)         & (3)          & (4)         & (5)        & (6)\\  
   \midrule 
   $Saving_{it}$            & 0.0726        & -0.0802     & 0.1260$^{*}$ & 0.0188      & 0.0364     & -0.0080\\   
                            & (0.0541)      & (0.1292)    & (0.0705)     & (0.0155)    & (0.1408)   & (0.0232)\\   
   $\log(Population_{it})$  & 0.4526$^{*}$  & -0.4247     & 0.0968       & -0.0257     & -0.4829    & -0.0015\\   
                            & (0.2711)      & (0.7100)    & (0.2283)     & (0.0620)    & (0.9378)   & (0.0780)\\   
   Log population control   & Yes           & Yes         & Yes          & Yes         & Yes        & Yes\\  
   Region FE                & Yes           & Yes         & Yes          & Yes         & Yes        & Yes\\  
   Year FE                  & Yes           & Yes         & Yes          & Yes         & Yes        & Yes\\  
   Regions                  & 160           & 160         & 160          & 151         & 160        & 151\\  
    \\
   Observations             & 3,040         & 2,880       & 3,040        & 2,730       & 2,880      & 2,587\\  
   Adjusted R$^2$           & 0.91769       & 0.06172     & 0.33381      & 0.04765     & 0.28506    & 0.05955\\  
   \bottomrule
\end{tabular}
\par\endgroup

                \end{table}

                \begin{table}[H]
                    \centering
                    \caption{All-Region Hospital Regression, Travel-Time Saving: 종합병원}
                    \input{tables/appendix_all_regions_saving_hospital_general.tex}
                \end{table}

                \begin{table}[H]
                    \centering
                    \caption{All-Region Hospital Regression, Travel-Time Saving: 상급종합병원}
                    \begingroup
\centering
\begin{tabular}{lcccccc}
   \toprule
                            & hospital\_no  & net\_entry  & exit\_no  & exit\_rate    & entry\_no  & entry\_rate \\    
                            & (1)           & (2)         & (3)       & (4)           & (5)        & (6)\\  
   \midrule 
   $Saving_{it}$            & 0.0219        & -0.0048     & 0.0010    & 0.0065        & -0.0036    & -0.1826\\   
                            & (0.0139)      & (0.0046)    & (0.0011)  & (0.0069)      & (0.0049)   & (0.1390)\\   
   $\log(Population_{it})$  & 0.0268        & -0.0011     & -0.0061   & -0.0328$^{*}$ & -0.0121    & -0.1291\\   
                            & (0.0689)      & (0.0134)    & (0.0061)  & (0.0169)      & (0.0227)   & (0.2889)\\   
   Log population control   & Yes           & Yes         & Yes       & Yes           & Yes        & Yes\\  
   Region FE                & Yes           & Yes         & Yes       & Yes           & Yes        & Yes\\  
   Year FE                  & Yes           & Yes         & Yes       & Yes           & Yes        & Yes\\  
   Regions                  & 160           & 160         & 160       & 20            & 160        & 20\\  
    \\
   Observations             & 3,040         & 2,880       & 3,040     & 279           & 2,880      & 266\\  
   Adjusted R$^2$           & 0.89446       & -0.02525    & 0.00013   & -0.01029      & -0.03110   & 0.06593\\  
   \bottomrule
\end{tabular}
\par\endgroup

                \end{table}

            \subsubsection{KTX Within Two Hours}
                \begin{table}[H]
                    \centering
                    \caption{All-Region Hospital Regression, KTX Within Two Hours: 의원}
                    \begingroup
\centering
\begin{tabular}{lcccccc}
   \toprule
                            & hospital\_no   & net\_entry  & exit\_no  & exit\_rate  & entry\_no     & entry\_rate \\    
                            & (1)            & (2)         & (3)       & (4)         & (5)           & (6)\\  
   \midrule 
   $KTXWithin2h_{it}$       & -0.0071        & 0.1729      & 0.2313    & 0.0033      & 0.4992$^{**}$ & 0.0004\\   
                            & (0.0265)       & (0.2411)    & (0.2986)  & (0.0047)    & (0.2365)      & (0.0085)\\   
   $\log(Population_{it})$  & 0.9245$^{***}$ & 1.197       & 0.6990    & 0.0133      & 3.083$^{***}$ & 0.0216$^{*}$\\   
                            & (0.1570)       & (0.8934)    & (1.211)   & (0.0081)    & (0.7722)      & (0.0129)\\   
   Log population control   & Yes            & Yes         & Yes       & Yes         & Yes           & Yes\\  
   Region FE                & Yes            & Yes         & Yes       & Yes         & Yes           & Yes\\  
   Year FE                  & Yes            & Yes         & Yes       & Yes         & Yes           & Yes\\  
   Regions                  & 160            & 160         & 160       & 159         & 160           & 159\\  
    \\
   Observations             & 3,040          & 2,880       & 3,040     & 3,021       & 2,880         & 2,862\\  
   Adjusted R$^2$           & 0.99451        & 0.33388     & 0.76234   & 0.14686     & 0.71000       & 0.06275\\  
   \bottomrule
\end{tabular}
\par\endgroup

                \end{table}

                \begin{table}[H]
                    \centering
                    \caption{All-Region Hospital Regression, KTX Within Two Hours: 병원}
                    \begingroup
\centering
\begin{tabular}{lcccccc}
   \toprule
                            & hospital\_no  & net\_entry  & exit\_no  & exit\_rate  & entry\_no  & entry\_rate \\    
                            & (1)           & (2)         & (3)       & (4)         & (5)        & (6)\\  
   \midrule 
   $KTXWithin2h_{it}$       & 0.1684        & 0.1504      & 0.0277    & 0.0135      & 0.1620     & -0.0104\\   
                            & (0.1147)      & (0.1799)    & (0.0460)  & (0.0247)    & (0.1943)   & (0.0540)\\   
   $\log(Population_{it})$  & 0.4697$^{*}$  & -0.3865     & 0.0885    & -0.0243     & -0.4580    & -0.0027\\   
                            & (0.2653)      & (0.7172)    & (0.2263)  & (0.0625)    & (0.9382)   & (0.0791)\\   
   Log population control   & Yes           & Yes         & Yes       & Yes         & Yes        & Yes\\  
   Region FE                & Yes           & Yes         & Yes       & Yes         & Yes        & Yes\\  
   Year FE                  & Yes           & Yes         & Yes       & Yes         & Yes        & Yes\\  
   Regions                  & 160           & 160         & 160       & 151         & 160        & 151\\  
    \\
   Observations             & 3,040         & 2,880       & 3,040     & 2,730       & 2,880      & 2,587\\  
   Adjusted R$^2$           & 0.91770       & 0.06161     & 0.33279   & 0.04735     & 0.28518    & 0.05952\\  
   \bottomrule
\end{tabular}
\par\endgroup

                \end{table}

                \begin{table}[H]
                    \centering
                    \caption{All-Region Hospital Regression, KTX Within Two Hours: 종합병원}
                    \begingroup
\centering
\begin{tabular}{lcccccc}
   \toprule
                            & hospital\_no  & net\_entry  & exit\_no  & exit\_rate  & entry\_no     & entry\_rate \\    
                            & (1)           & (2)         & (3)       & (4)         & (5)           & (6)\\  
   \midrule 
   $KTXWithin2h_{it}$       & -0.0187       & -0.0141     & -0.0186   & -0.0359     & -0.0321$^{*}$ & -0.0492\\   
                            & (0.0377)      & (0.0130)    & (0.0136)  & (0.0273)    & (0.0186)      & (0.0358)\\   
   $\log(Population_{it})$  & 0.3359$^{**}$ & -0.0184     & -0.0030   & 0.0070      & 0.0061        & -0.1550\\   
                            & (0.1451)      & (0.0672)    & (0.0304)  & (0.0450)    & (0.0859)      & (0.1102)\\   
   Log population control   & Yes           & Yes         & Yes       & Yes         & Yes           & Yes\\  
   Region FE                & Yes           & Yes         & Yes       & Yes         & Yes           & Yes\\  
   Year FE                  & Yes           & Yes         & Yes       & Yes         & Yes           & Yes\\  
   Regions                  & 160           & 160         & 160       & 99          & 160           & 99\\  
    \\
   Observations             & 3,040         & 2,880       & 3,040     & 1,678       & 2,880         & 1,594\\  
   Adjusted R$^2$           & 0.92593       & -0.00751    & 0.10421   & 0.05870     & 0.06413       & 0.05422\\  
   \bottomrule
\end{tabular}
\par\endgroup

                \end{table}

                \begin{table}[H]
                    \centering
                    \caption{All-Region Hospital Regression, KTX Within Two Hours: 상급종합병원}
                    \begingroup
\centering
\begin{tabular}{lcccccc}
   \toprule
                            & hospital\_no  & net\_entry  & exit\_no  & exit\_rate     & entry\_no  & entry\_rate \\    
                            & (1)           & (2)         & (3)       & (4)            & (5)        & (6)\\  
   \midrule 
   $KTXWithin2h_{it}$       & 0.0055        & 0.0060      & 0.0293    & 0.1661$^{***}$ & 0.0360     & 0.2435$^{***}$\\   
                            & (0.0088)      & (0.0058)    & (0.0290)  & (0.0004)       & (0.0353)   & (0.0415)\\   
   $\log(Population_{it})$  & 0.0255        & 0.0007      & -0.0020   & -0.0051$^{*}$  & -0.0050    & -0.0090\\   
                            & (0.0704)      & (0.0125)    & (0.0025)  & (0.0028)       & (0.0167)   & (0.2537)\\   
   Log population control   & Yes           & Yes         & Yes       & Yes            & Yes        & Yes\\  
   Region FE                & Yes           & Yes         & Yes       & Yes            & Yes        & Yes\\  
   Year FE                  & Yes           & Yes         & Yes       & Yes            & Yes        & Yes\\  
   Regions                  & 160           & 160         & 160       & 20             & 160        & 20\\  
    \\
   Observations             & 3,040         & 2,880       & 3,040     & 279            & 2,880      & 266\\  
   Adjusted R$^2$           & 0.89408       & -0.02539    & 0.02544   & 0.28181        & -0.02904   & 0.06451\\  
   \bottomrule
\end{tabular}
\par\endgroup

                \end{table}

        \subsection{Per-Capita Outcomes}
            이 subsection은 본문 section 1의 per-capita 분석을 병원 회귀식에 대해 160개 지역
            전체 샘플로 반복한 것이다. 종속변수는 병원 수, 순진입, 폐업 수, 진입 수를 인구
            만 명당 값으로 바꾼 변수이고, 폐업률과 진입률은 원래 rate 변수를 그대로 사용하였다.
            이 분석에서는 종속변수 자체가 인구로 scale되어 있으므로
            $\log(Population_{it})$ control은 포함하지 않는다.

            \begin{align*}
                Y^m_{it}
                =
                \alpha_i + \tau_t
                + \beta Access_{it}
                + \epsilon_{it}
            \end{align*}

            \subsubsection{Travel Time}
                \begin{table}[H]
                    \centering
                    \caption{All-Region Per-Capita Hospital Regression, Travel Time: 의원}
                    \begingroup
\centering
\begin{tabular}{lcccccc}
   \toprule
                          & hospital\_no  & net\_entry  & exit\_no  & exit\_rate  & entry\_no  & entry\_rate \\    
                          & (1)           & (2)         & (3)       & (4)         & (5)        & (6)\\  
   \midrule 
   $TravelTime_{it}$      & 0.1416        & 0.0099      & 0.0022    & -0.0011     & 0.0105     & 0.0004\\   
                          & (0.0856)      & (0.0143)    & (0.0137)  & (0.0028)    & (0.0177)   & (0.0039)\\   
   Log population control & No            & No          & No        & No          & No         & No\\  
   Region FE              & Yes           & Yes         & Yes       & Yes         & Yes        & Yes\\  
   Year FE                & Yes           & Yes         & Yes       & Yes         & Yes        & Yes\\  
   Regions                & 160           & 160         & 160       & 159         & 160        & 159\\  
    \\
   Observations           & 3,040         & 2,880       & 3,040     & 3,021       & 2,880      & 2,862\\  
   Adjusted R$^2$         & 0.96475       & 0.11628     & 0.34113   & 0.14672     & 0.26162    & 0.06264\\  
   \bottomrule
\end{tabular}
\par\endgroup

                \end{table}

                \begin{table}[H]
                    \centering
                    \caption{All-Region Per-Capita Hospital Regression, Travel Time: 병원}
                    \begingroup
\centering
\begin{tabular}{lcccccc}
   \toprule
                          & hospital\_no  & net\_entry  & exit\_no       & exit\_rate  & entry\_no     & entry\_rate \\    
                          & (1)           & (2)         & (3)            & (4)         & (5)           & (6)\\  
   \midrule 
   $TravelTime_{it}$      & -0.0244       & -0.0005     & -0.0093$^{**}$ & -0.0188     & -0.0099$^{*}$ & 0.0080\\   
                          & (0.0209)      & (0.0055)    & (0.0044)       & (0.0155)    & (0.0057)      & (0.0232)\\   
   Log population control & No            & No          & No             & No          & No            & No\\  
   Region FE              & Yes           & Yes         & Yes            & Yes         & Yes           & Yes\\  
   Year FE                & Yes           & Yes         & Yes            & Yes         & Yes           & Yes\\  
   Regions                & 160           & 160         & 160            & 151         & 160           & 151\\  
    \\
   Observations           & 3,040         & 2,880       & 3,040          & 2,730       & 2,880         & 2,587\\  
   Adjusted R$^2$         & 0.75744       & 0.00577     & 0.09879        & 0.04795     & 0.06199       & 0.05994\\  
   \bottomrule
\end{tabular}
\par\endgroup

                \end{table}

                \begin{table}[H]
                    \centering
                    \caption{All-Region Per-Capita Hospital Regression, Travel Time: 종합병원}
                    \input{tables/appendix_all_regions_pc_travel_time_hospital_general.tex}
                \end{table}

                \begin{table}[H]
                    \centering
                    \caption{All-Region Per-Capita Hospital Regression, Travel Time: 상급종합병원}
                    \input{tables/appendix_all_regions_pc_travel_time_hospital_tertiary.tex}
                \end{table}

            \subsubsection{Travel-Time Saving}
                \begin{table}[H]
                    \centering
                    \caption{All-Region Per-Capita Hospital Regression, Travel-Time Saving: 의원}
                    \input{tables/appendix_all_regions_pc_saving_hospital_clinic.tex}
                \end{table}

                \begin{table}[H]
                    \centering
                    \caption{All-Region Per-Capita Hospital Regression, Travel-Time Saving: 병원}
                    \input{tables/appendix_all_regions_pc_saving_hospital_secondary.tex}
                \end{table}

                \begin{table}[H]
                    \centering
                    \caption{All-Region Per-Capita Hospital Regression, Travel-Time Saving: 종합병원}
                    \begingroup
\centering
\begin{tabular}{lcccccc}
   \toprule
                          & hospital\_no  & net\_entry  & exit\_no  & exit\_rate  & entry\_no  & entry\_rate \\    
                          & (1)           & (2)         & (3)       & (4)         & (5)        & (6)\\  
   \midrule 
   $Saving_{it}$          & 0.0040        & -0.0006     & -0.0018   & -0.0172     & -0.0023    & -0.0187\\   
                          & (0.0048)      & (0.0013)    & (0.0011)  & (0.0108)    & (0.0017)   & (0.0205)\\   
   Log population control & No            & No          & No        & No          & No         & No\\  
   Region FE              & Yes           & Yes         & Yes       & Yes         & Yes        & Yes\\  
   Year FE                & Yes           & Yes         & Yes       & Yes         & Yes        & Yes\\  
   Regions                & 160           & 160         & 160       & 99          & 160        & 99\\  
    \\
   Observations           & 3,040         & 2,880       & 3,040     & 1,678       & 2,880      & 1,594\\  
   Adjusted R$^2$         & 0.88034       & -0.00474    & 0.03402   & 0.05975     & 0.01457    & 0.05291\\  
   \bottomrule
\end{tabular}
\par\endgroup

                \end{table}

                \begin{table}[H]
                    \centering
                    \caption{All-Region Per-Capita Hospital Regression, Travel-Time Saving: 상급종합병원}
                    \begingroup
\centering
\begin{tabular}{lcccccc}
   \toprule
                          & hospital\_no  & net\_entry  & exit\_no                & exit\_rate  & entry\_no  & entry\_rate \\    
                          & (1)           & (2)         & (3)                     & (4)         & (5)        & (6)\\  
   \midrule 
   $Saving_{it}$          & 0.0033        & -0.0004     & $1.03\times 10^{-5}$    & 0.0074      & -0.0004    & -0.1821\\   
                          & (0.0022)      & (0.0005)    & ($1.04\times 10^{-5}$)  & (0.0083)    & (0.0005)   & (0.1390)\\   
   Log population control & No            & No          & No                      & No          & No         & No\\  
   Region FE              & Yes           & Yes         & Yes                     & Yes         & Yes        & Yes\\  
   Year FE                & Yes           & Yes         & Yes                     & Yes         & Yes        & Yes\\  
   Regions                & 160           & 160         & 160                     & 20          & 160        & 20\\  
    \\
   Observations           & 3,040         & 2,880       & 3,040                   & 279         & 2,880      & 266\\  
   Adjusted R$^2$         & 0.92934       & -0.02276    & -0.00016                & -0.01386    & -0.02289   & 0.06863\\  
   \bottomrule
\end{tabular}
\par\endgroup

                \end{table}

            \subsubsection{KTX Within Two Hours}
                \begin{table}[H]
                    \centering
                    \caption{All-Region Per-Capita Hospital Regression, KTX Within Two Hours: 의원}
                    \begingroup
\centering
\begin{tabular}{lcccccc}
   \toprule
                          & hospital\_no  & net\_entry  & exit\_no  & exit\_rate  & entry\_no  & entry\_rate \\    
                          & (1)           & (2)         & (3)       & (4)         & (5)        & (6)\\  
   \midrule 
   $KTXWithin2h_{it}$     & -0.1307       & 0.0050      & 0.0126    & 0.0022      & 0.0049     & -0.0015\\   
                          & (0.1023)      & (0.0223)    & (0.0307)  & (0.0049)    & (0.0451)   & (0.0086)\\   
   Log population control & No            & No          & No        & No          & No         & No\\  
   Region FE              & Yes           & Yes         & Yes       & Yes         & Yes        & Yes\\  
   Year FE                & Yes           & Yes         & Yes       & Yes         & Yes        & Yes\\  
   Regions                & 160           & 160         & 160       & 159         & 160        & 159\\  
    \\
   Observations           & 3,040         & 2,880       & 3,040     & 3,021       & 2,880      & 2,862\\  
   Adjusted R$^2$         & 0.96463       & 0.11619     & 0.34116   & 0.14671     & 0.26156    & 0.06265\\  
   \bottomrule
\end{tabular}
\par\endgroup

                \end{table}

                \begin{table}[H]
                    \centering
                    \caption{All-Region Per-Capita Hospital Regression, KTX Within Two Hours: 병원}
                    \begingroup
\centering
\begin{tabular}{lcccccc}
   \toprule
                          & hospital\_no  & net\_entry  & exit\_no  & exit\_rate  & entry\_no  & entry\_rate \\    
                          & (1)           & (2)         & (3)       & (4)         & (5)        & (6)\\  
   \midrule 
   $KTXWithin2h_{it}$     & 0.0797$^{**}$ & 0.0022      & 0.0022    & 0.0152      & 0.0024     & -0.0102\\   
                          & (0.0312)      & (0.0098)    & (0.0077)  & (0.0235)    & (0.0066)   & (0.0535)\\   
   Log population control & No            & No          & No        & No          & No         & No\\  
   Region FE              & Yes           & Yes         & Yes       & Yes         & Yes        & Yes\\  
   Year FE                & Yes           & Yes         & Yes       & Yes         & Yes        & Yes\\  
   Regions                & 160           & 160         & 160       & 151         & 160        & 151\\  
    \\
   Observations           & 3,040         & 2,880       & 3,040     & 2,730       & 2,880      & 2,587\\  
   Adjusted R$^2$         & 0.75804       & 0.00578     & 0.09793   & 0.04765     & 0.06146    & 0.05991\\  
   \bottomrule
\end{tabular}
\par\endgroup

                \end{table}

                \begin{table}[H]
                    \centering
                    \caption{All-Region Per-Capita Hospital Regression, KTX Within Two Hours: 종합병원}
                    \begingroup
\centering
\begin{tabular}{lcccccc}
   \toprule
                          & hospital\_no  & net\_entry  & exit\_no  & exit\_rate  & entry\_no  & entry\_rate \\    
                          & (1)           & (2)         & (3)       & (4)         & (5)        & (6)\\  
   \midrule 
   $KTXWithin2h_{it}$     & 0.0072        & -0.0004     & -0.0029   & -0.0367     & -0.0036    & -0.0299\\   
                          & (0.0056)      & (0.0004)    & (0.0027)  & (0.0273)    & (0.0029)   & (0.0305)\\   
   Log population control & No            & No          & No        & No          & No         & No\\  
   Region FE              & Yes           & Yes         & Yes       & Yes         & Yes        & Yes\\  
   Year FE                & Yes           & Yes         & Yes       & Yes         & Yes        & Yes\\  
   Regions                & 160           & 160         & 160       & 99          & 160        & 99\\  
    \\
   Observations           & 3,040         & 2,880       & 3,040     & 1,678       & 2,880      & 1,594\\  
   Adjusted R$^2$         & 0.88029       & -0.00481    & 0.03358   & 0.05929     & 0.01425    & 0.05253\\  
   \bottomrule
\end{tabular}
\par\endgroup

                \end{table}

                \begin{table}[H]
                    \centering
                    \caption{All-Region Per-Capita Hospital Regression, KTX Within Two Hours: 상급종합병원}
                    \begingroup
\centering
\begin{tabular}{lcccccc}
   \toprule
                          & hospital\_no  & net\_entry              & exit\_no  & exit\_rate              & entry\_no  & entry\_rate \\    
                          & (1)           & (2)                     & (3)       & (4)                     & (5)        & (6)\\  
   \midrule 
   $KTXWithin2h_{it}$     & 0.0004        & $7.16\times 10^{-5}$    & 0.0003    & 0.1667$^{***}$          & 0.0004     & 0.2450$^{***}$\\   
                          & (0.0004)      & ($7.68\times 10^{-5}$)  & (0.0003)  & ($3.44\times 10^{-9}$)  & (0.0003)   & (0.0232)\\   
   Log population control & No            & No                      & No        & No                      & No         & No\\  
   Region FE              & Yes           & Yes                     & Yes       & Yes                     & Yes        & Yes\\  
   Year FE                & Yes           & Yes                     & Yes       & Yes                     & Yes        & Yes\\  
   Regions                & 160           & 160                     & 160       & 20                      & 160        & 20\\  
    \\
   Observations           & 3,040         & 2,880                   & 3,040     & 279                     & 2,880      & 266\\  
   Adjusted R$^2$         & 0.92885       & -0.02312                & 0.02572   & 0.28462                 & -0.02318   & 0.06861\\  
   \bottomrule
\end{tabular}
\par\endgroup

                \end{table}

        \subsection{Difference From 2007}
            이 subsection은 본문 section 2의 2007년 대비 변화분 분석을 병원 stock 변수에 대해
            160개 지역 전체 샘플로 반복한 것이다. 종속변수는 각 병원 type의 병원 수에서
            2007년 병원 수를 뺀 값이며, 독립변수는 $Saving_{it}$이다. Control로는
            $\log(Population_{it}) - \log(Population_{i,2007})$를 포함한다.

            \begin{align*}
                HospitalNo^m_{it} - HospitalNo^m_{i,2007}
                =
                \alpha_i + \tau_t
                + \beta Saving_{it}
                + \gamma \left[
                    \log(Population_{it}) - \log(Population_{i,2007})
                \right]
                + \epsilon_{it}
            \end{align*}

            \begin{table}[H]
                \centering
                \caption{All-Region Difference From 2007, Travel-Time Saving: 병원 수}
                \begingroup
\centering
\begin{tabular}{lcccc}
   \toprule
                                     & 의원    & 병원     & 종합병원 & 상급종합병원 \\   
                                     & (1)     & (2)      & (3)      & (4)\\  
   \midrule 
   $Saving_{it}$                     & -1.610  & 0.3526   & 0.0666   & 0.0342$^{*}$\\   
                                     & (2.139) & (0.3423) & (0.0682) & (0.0204)\\   
   Log population difference control & Yes     & Yes      & Yes      & Yes\\  
   Region FE                         & Yes     & Yes      & Yes      & Yes\\  
   Year FE                           & Yes     & Yes      & Yes      & Yes\\  
   Regions                           & 160     & 160      & 160      & 160\\  
    \\
   Observations                      & 3,040   & 3,040    & 3,040    & 3,040\\  
   Adjusted R$^2$                    & 0.74036 & 0.75776  & 0.68686  & 0.47875\\  
   \bottomrule
\end{tabular}
\par\endgroup

            \end{table}

\end{document}
